\cleardoublepage
\chapter{Środowisko badawcze i implementacja}
\label{cha:SrodowiskoImplementacja}

W niniejszym rozdziale zaprezentowano środowisko sprzętowe i~programowe wykorzystane do~przeprowadzenia eksperymentów. 
Omówiono zastosowane technologie, sposób implementacji analizowanych algorytmów oraz charakterystykę przygotowanych danych testowych.

\section{Opis środowiska}
\label{sec:Srodowisko}

Wszystkie eksperymenty zostały przeprowadzone na komputerze o następującej konfiguracji:

\begin{itemize}
    \item \textbf{Procesor:} Intel(R) Core(TM) i9-9980XE CPU @ 3.00\,GHz
    \item \textbf{Rdzenie fizyczne:} 18
    \item \textbf{Rdzenie logiczne:} 36 
    \item \textbf{Taktowanie bazowe CPU:} 3000\,MHz
    \item \textbf{Pamięć operacyjna RAM:} 127,68\,GB
    \item \textbf{System operacyjny:} Windows Server 2019
    \item \textbf{Architektura:} AMD64 (x86-64)
    \item \textbf{Interpreter Python:} 3.14.7
\end{itemize}

Środowisko to umożliwia przeprowadzenie badań algorytmów sortowania równoległego dla różnych poziomów równoległości. 
Procesor wyposażony w 18 rdzeni fizycznych i 36 procesorów logicznych pozwala na analizę zmian wydajności wraz ze wzrostem 
liczby wykorzystywanych jednostek wykonawczych.

Dane dotyczące sprzętu i~oprogramowania zostały automatycznie zebrane i~zapisywane w bazie danych SQLite przed rozpoczęciem eksperymentów.

\section{Wybór języka i bibliotek}
\label{sec:JezykBiblioteki}

Do~implementacji algorytmów sortowania równoległego, generowania danych, przeprowadzenia eksperymentów oraz wizualizacji wyników wybrano 
język programowania Python w wersji 3.14.7~\cite{bib:Python}. Decyzja ta wynikała z~kilku powodów. Python jest jednym z~najpopularniejszych 
języków, charakteryzuje się przejrzystą i~elegancką składnią, co ułatwia tworzenie, analizę i~utrzymanie kodu. 
Posiada bogatą bibliotekę standardową oraz rozbudowany zbiór dostępnych bibliotek, pakietów i~narzędzi programistycznych rozszerzających 
możliwości języka, które znacząco przyspieszają rozwój oprogramowania. Dodatkowo język ten jest w~pełni przenośny między systemami 
operacyjnymi, takimi jak Windows, Linux, macOS, co ma znaczenie przy planowaniu ewentualnych dalszych testów na innych platformach.

Niezwykle istotną rolę w projekcie odgrywa biblioteka \texttt{multiprocessing}~\cite{bib:Multiprocessing}, stanowiąca część standardowej 
biblioteki Pythona. Pakiet ten przy użyciu interfejsu programistycznego podobnego do~pythonowego modułu \texttt{threading} umożliwia 
tworzenie procesów oraz realizację obliczeń równoległych z pominięciem ograniczeń wynikających z globalnej blokady interpretera 
(ang. \texttt{global interpreter lock}, GIL). Dzięki temu możliwe jest efektywne wykorzystanie wielu rdzeni procesora. 
Do współdzielenia danych wykorzystano moduł \texttt{shared\_memory} z~pakietu \texttt{multiprocessing}~\cite{bib:SharedMemory}, 
który pozwala na bezpośredni dostęp wielu procesów do~wspólnego obszaru pamięci. 
Mechanizm ten eliminuje konieczność serializacji oraz przesyłania dużych struktur danych między procesami, 
ograniczając związany z tym narzut, co ma szczególne znaczenie podczas sortowania dużych zbiorów danych.

Biblioteka \texttt{NumPy}~\cite{bib:NumPy} obsługująca dane numeryczne i~szybkie operacje na tablicach, została wykorzystana w~optymalizacji 
algorytmów Parallel Bucket Sort oraz Parallel Sample Sort. Posłużyła w nich do~efektywnego rozdzielania elementów do~kubełków, 
co znacznie usprawniło etap dystrybucji danych. Ponadto \texttt{NumPy} wykorzystano w module analizy złożoności obliczeniowej oraz 
przy generowaniu wykresów i~rankingów. 

Wyniki eksperymentów były przechowywane i~przetwarzane z użyciem biblioteki \texttt{pandas}~\cite{bib:Pandas}, która oferuje wygodne 
struktury danych (DataFrame) oraz narzędzia do~filtrowania, agregacji i~analizy tabelarycznej. Do~generowania wykresów i~wizualizacji 
wyników wykorzystano bibliotekę \texttt{matplotlib}~\cite{bib:Matplotlib}.

Do monitorowania wykorzystanych zasobów systemowych wykorzystano bibliotekę \texttt{psutil}~\cite{bib:Psutil} oraz 
narzędzia \texttt{py-cpuinfo} i~\texttt{py-spy}~\cite{bib:PySpy}. 

Całość kodu źródłowego została opracowana w środowisku programistycznym JetBrains PyCharm~\cite{bib:PyCharm}, 
które zapewnia wsparcie dla języka Python, debugowanie oraz wygodne zarządzanie projektem.

Do przechowywania wygenerowanych danych testowych oraz wyników eksperymentów wykorzystano bazę danych SQLite~\cite{bib:SQLite}. 
Jest to wbudowany silnik bazodanowy, który nie wymaga osobnego serwera, ponieważ odczytuje i~zapisuje dane bezpośrednio do plików dyskowych. 
Rozwiązanie to upraszcza zarządzanie danymi i~zapewnia łatwą przenośność między środowiskami.

Wybór powyższego zestawu technologii umożliwia spójną implementację algorytmów, precyzyjny pomiar ich charakterystyk wydajnościowych 
oraz wygodną analizę i~prezentację wyników badań.

\section{Implementacja wybranych algorytmów}
\label{sec:ImplementacjaAlgorytmow}

Wszystkie algorytmy zaimplementowano w języku Python z wykorzystaniem mechanizmu wieloprocesowego oraz pamięci współdzielonej. 
Każdy algorytm występuje w dwóch wariantach: sekwencyjnym, który stanowi punkt odniesienia, oraz równoległym. 
Wspólną cechą implementacji równoległych jest przechowywanie sortowanych danych w bloku pamięci współdzielonej, 
mapowanym na tablicę typów \texttt{ctypes} (\texttt{c\_longlong} dla liczb całkowitych oraz \texttt{c\_double} dla liczb zmiennoprzecinkowych). 
Dzięki temu procesy potomne mogą odczytywać i modyfikować te same dane bez konieczności ich kopiowania oraz przesyłania pomiędzy procesami.

\subsection*{Ogólna architektura rozwiązania}

Przy algorytmach opartych na strategii ``dziel i zwyciężaj'' czyli Parallel Quick Sort oraz Parallel Merge Sort, 
zastosowano rekurencyjne tworzenie procesów z ograniczeniem głębokości równoległości za pomocą parametru \texttt{max\_depth}. 
Dodatkowo wprowadzono progi minimalnego rozmiaru fragmentu danych zawarte w słowniku\\ \texttt{PARALLEL\_CUTOFF}, 
poniżej których dalsze tworzenie procesów jest nieopłacalne i~następuje przejście do sortowania sekwencyjnego. 
Takie rozwiązanie ogranicza narzut związany z tworzeniem zbyt dużej liczby procesów przy małych fragmentach danych.

Dla algorytmów opartych na kubełkach czyli Parallel Bucket Sort oraz Parallel Sample Sort, zastosowano model, 
w którym liczba procesów jest określana na podstawie liczby jednostek wykonawczych wykorzystywanych w danym eksperymencie. 
Dane są najpierw rozdzielane do kubełków, a~następnie grupy kubełków są przydzielane poszczególnym procesom. 
Również tutaj wprowadzono progi minimalnego rozmiaru zawarte w \texttt{PARALLEL\_SIZE\_CUTOFF} 
oraz \texttt{GROUP\_SIZE\_CUTOFF}, ograniczające wykorzystanie przetwarzania równoległego dla niewielkich zbiorów lub grup danych.

Wartości progów zastosowanych w~poszczególnych implementacjach przedstawiono w~tabeli~\ref{tab:parallel-cutoffs}.

Kolumna Quick/Merge zawiera wartości parametru \texttt{PARALLEL\_CUTOFF}, 
natomiast kolumna Bucket/Sample wartości parametru \texttt{PARALLEL\_SIZE\_CUTOFF}. 
Dwie ostatnie kolumny przedstawiają wartości parametru \texttt{GROUP\_SIZE\_CUTOFF}, różne dla implementacji Bucket Sort i~Sample Sort.

\input{tables/chapter4/parallel_cutoffs}

Próg \texttt{PARALLEL\_CUTOFF} określa minimalny rozmiar fragmentu danych,
dla którego w~implementacjach Quick Sort i~Merge Sort dopuszczalne jest dalsze tworzenie procesów.
Po osiągnięciu wartości progowej dalsze sortowanie danego fragmentu wykonywane jest sekwencyjnie.

Parametr \texttt{PARALLEL\_SIZE\_CUTOFF} wyznacza próg rozmiaru danych
brany pod uwagę przy wyborze pomiędzy wykonaniem równoległym a~sekwencyjnym w~implementacjach Bucket Sort oraz Sample Sort. 
Natomiast \texttt{GROUP\_SIZE\_CUTOFF} ogranicza tworzenie osobnych procesów dla grup danych o~małym rozmiarze. 
Wartości tego parametru są różne dla implementacji Bucket Sort i~Sample Sort.

\subsection*{Parallel Quick Sort}

Wariant sekwencyjny realizuje klasyczny Quick Sort z wyborem pivota jako elementu środkowego.

\lstinputlisting[
  language=Python, caption={Funkcja partycjonowania w Quick Sort}, label={lst:qs-partition}
]{listings/chapter4/quicksort_partition.py}

W wersji równoległej po wykonaniu partycjonowania lewa i prawa część tablicy mogą być sortowane niezależnie przez osobne procesy. 
Każdy proces dołącza się do~istniejącego bloku pamięci współdzielonej, sortuje przypisany dla siebie zakres, a~następnie kończy działanie. 
Gdy rozmiar zakresu spada poniżej ustalonego progu lub osiągnięta zostaje maksymalna głębokość równoległości, 
dalsze sortowanie odbywa się sekwencyjnie.

\lstinputlisting[
  language=Python, caption={Tworzenie procesów w Parallel Quick Sort}, label={lst:qs-parallel-spawn}
]{listings/chapter4/quicksort_parallel_spawn.py}

\subsection*{Parallel Merge Sort}

Wariant sekwencyjny realizuje klasyczny Merge Sort z rekurencyjnym podziałem tablicy na połowy oraz scalaniem posortowanych części.

\lstinputlisting[
  language=Python, caption={Funkcja scalania w Merge Sort}, label={lst:ms-merge}
]{listings/chapter4/mergesort_merge.py}

W wersji równoległej obie połowy mogą być sortowane jednocześnie przez osobne procesy. 
Po zakończeniu obu procesów następuje etap scalania, który w~danej implementacji wykonywany jest sekwencyjnie. 
Głębokość równoległości oraz~minimalny rozmiar zakresu ograniczają nadmierne tworzenie procesów.

\lstinputlisting[
  language=Python, caption={Równoległe sortowanie połówek i sekwencyjne scalanie w Parallel Merge Sort}, label={lst:ms-parallel-spawn}
]{listings/chapter4/mergesort_parallel_spawn.py}

\subsection*{Parallel Bucket Sort}

Wariant sekwencyjny rozdziela elementy do liczby kubełków wyznaczanej na podstawie rozmiaru danych, a~następnie 
sortuje każdy kubełek z wykorzystaniem zaimplementowanego algorytmu Merge Sort. Na sam koniec łączy wyniki w jedną posortowaną tablicę.

\lstinputlisting[
  language=Python, caption={Rozdzielanie elementów do kubełków w Bucket Sort}, label={lst:bs-distribute}
]{listings/chapter4/bucketsort_distribute.py}

W wersji równoległej liczba utworzonych kubełków jest wyliczana na podstawie rozmiaru danych oraz liczby jednostek wykonawczych. 
Po utworzeniu kubełków ich zawartość jest łączona w jedną tablicę, która umieszczana zostaje w pamięci współdzielonej, 
a zakresy poszczególnych kubełków są grupowane i przydzielane procesom. Procesy sortują przypisane im grupy kubełków niezależnie. 
Dla małych grup sortowanie wykonywane jest sekwencyjnie, co ogranicza narzut związany z~tworzeniem procesów.
Do efektywnego rozdzielania elementów wykorzystano operacje wektorowe biblioteki \texttt{NumPy}.

\lstinputlisting[
  language=Python, caption={Przydzielanie grup kubełków procesom w Parallel Bucket Sort}, label={lst:bs-parallel-spawn}
]{listings/chapter4/bucketsort_parallel_spawn.py}

\subsection*{Parallel Sample Sort}

Wariant sekwencyjny działa w kilku etapach: podział danych na części, lokalne sortowanie, wybór próbek, wyznaczenie wartości rozdzielających, 
rozdzielenie elementów do kubełków według wyznaczonych wartości oraz ostateczne sortowanie kubełków.

\lstinputlisting[
  language=Python, caption={Rozdzielanie elementów do kubełków na podstawie wartości rozdzielających w Sample Sort}, label={lst:ss-distribute}
]{listings/chapter4/samplesort_distribute.py}

W wersji równoległej etapy te realizowane są z wykorzystaniem wielu procesów. 
Na początku procesy równolegle sortują lokalne fragmenty danych i zbierają próbki. 
Na podstawie próbek wyznaczane są wartości rozdzielające, po czym dane elementy są rozdzielane do odpowiednich kubełków. 
Na końcu grupy kubełków są ponownie sortowane równolegle. 
W przypadku niewielkich grup danych stosowany jest wariant sekwencyjny, co ogranicza koszt związany z tworzeniem procesów.
Do dystrybucji elementów wykorzystano operacje wektorowe biblioteki \texttt{NumPy}.

\lstinputlisting[
  language=Python, caption={Równoległe sortowanie grup kubełków w Parallel Sample Sort}, label={lst:ss-parallel-spawn}
]{listings/chapter4/samplesort_parallel_spawn.py}

\section{Charakterystyka danych testowych}
\label{sec:DaneTestowane}

Do przeprowadzenia testów przygotowano zestawy danych o zróżnicowanej charakterystyce, 
tak aby możliwe było ocenienie zachowania algorytmów w różnych scenariuszach.

\subsection*{Typy i rozmiary danych}

Testy przeprowadzono dla dwóch typów wartości:
\begin{itemize}
    \item liczb całkowitych -- \texttt{int}
    \item liczb zmiennoprzecinkowych -- \texttt{float}
\end{itemize}

Dla każdego typu wygenerowano zbiory o następujących rozmiarach:
\begin{itemize}
    \item 1~000 elementów
    \item 10~000 elementów
    \item 100~000 elementów
    \item 1~000~000 elementów
\end{itemize}

Wygenerowane wartości znajdują się w zakresie od 0 do 1~000~000.

\subsection*{Charakterystyka zbiorów}

Przygotowano trzy główne kategorie danych:

\begin{itemize}
    \item \textbf{Dane losowe} -- elementy generowane w sposób całkowicie losowy     z użyciem generatora liczb pseudolosowych.
    
    \item \textbf{Dane z duplikatami} -- zbiory, w których około 40\% elementów stanowią powtórzenia. 
    Najpierw wygenerowano 60\% unikalnych wartości zbioru o~danym rozmiarze, 
    a następnie na ich podstawie dodano brakującą część poprzez losowe powtórzenia.
    
    \item \textbf{Dane częściowo posortowane} -- zbiory zawierające odpowiednio 20\%, 40\%, 60\% oraz 80\% elementów uporządkowanych. 
    Posortowane fragmenty przeplatano fragmentami losowymi, co pozwoliło uzyskać zbiory o określonym stopniu częściowego uporządkowania.
\end{itemize}

\lstinputlisting[
  language=Python, caption={Generowanie danych częściowo posortowanych}, label={lst:data-part-sorted}
]{listings/chapter4/generate_part_sorted.py}

Do testów przygotowano dwanaście wariantów zbiorów (2 typy danych $\times$ 6 charakterystyk).

\subsection*{Sposób przechowywania}

Wszystkie dane testowe zapisano w bazie danych SQLite. Każdy wariant przechowywany był w osobnej tabeli o strukturze:

\begin{itemize}
    \item \texttt{id} -- unikalny identyfikator rekordu
    \item \texttt{value} -- wartość elementu \texttt{INTEGER} lub \texttt{REAL}
    \item \texttt{set\_size} -- rozmiar zbioru, do którego należy dany element, ograniczony do wartości 1~000, 10~000, 100~000, 1~000~000
\end{itemize}